\documentclass[11pt,a4paper]{article}
\usepackage{amsmath,amssymb,booktabs,geometry,hyperref}
\geometry{margin=1in}
\title{\textbf{The Virial Partition from Atoms to the Horizon}}
\author{Heath W.\ Mahaffey\\ \small Independent Researcher, Entiat, WA, USA\\ \small \url{https://github.com/hmahaffeyges/IAM-Validation}}
\date{October 2026}
\begin{document}
\maketitle
\begin{abstract}
For any system bound by a potential that falls as $1/r$, the virial theorem fixes the time-averaged kinetic energy at exactly half the magnitude of the
potential energy. We follow this one half through every scale at which it can be checked: the hydrogen atom, atoms and molecules at the Hartree--Fock limit,
the Sun, white dwarfs, clusters of galaxies, the horizon of a black hole, and the horizon of the universe. We state the thermodynamic content of the kinetic
half: when a $1/r$ system binds from rest, the energy it releases equals its kinetic half, and Landauer's principle makes that the minimum cost of the
irreversible binding. For a Schwarzschild black hole the Bekenstein--Hawking entropy counted in bits, priced at $k_BT_H\ln2$ each, equals exactly
$\tfrac12Mc^2$ (the Smarr relation). At the cosmic horizon the same half fixes the coupling $\beta_m=\Omega_m/2$ of the Informational Actualization Model,
which enters every Planck chain at its predicted value and is consistent with Planck 2018 ($\Delta\chi^2=+0.54$ against $\Lambda$CDM, no added parameter).
Every number is recomputed by the script \texttt{docs/verification/scripts/verify\_virial\_atoms\_to\_horizon.py}.
\end{abstract}

\section{The theorem}
For a bound system with potential energy homogeneous of degree $k$ in the coordinates, the time average of $d(\sum_i\mathbf r_i\cdot\mathbf p_i)/dt$ vanishes,
and Euler's theorem gives $2\langle K\rangle=k\langle V\rangle$. Gravity and the Coulomb interaction have $k=-1$:
\begin{equation}
2\langle K\rangle+\langle V\rangle=0,\qquad \langle K\rangle=\tfrac12|\langle V\rangle|,\qquad E=\langle K\rangle+\langle V\rangle=-\langle K\rangle .
\label{eq:virial}
\end{equation}
The degree is fixed by geometry: in three spatial dimensions Gauss's law makes the field of a point source fall as $1/r^2$, so the potential falls as $1/r$.
The half does not depend on the number of bodies, their masses or the scale.

\section{Atoms and molecules}
\textbf{Hydrogen.} The ground state has $E_1=-13.6057$\,eV, $\langle K\rangle=13.6057$\,eV and $\langle V\rangle=-27.2114$\,eV: $\langle K\rangle/|\langle V\rangle|=1/2$.
An electron captured from rest into this state emits a photon of $13.6057$\,eV, equal to its kinetic half.

\textbf{Many-electron atoms and molecules.} For the 20 atoms H--Xe and 10 molecules (H$_2$ to CO$_2$) at the Hartree--Fock limit (Clementi \& Roetti 1974;
Bunge et al.\ 1993), $T/|V|=1/2$ to the precision of the tabulated energies. For a converged variational wavefunction this is exact: optimising under a uniform
scaling of the coordinates enforces the quantum virial theorem. These systems show that quantum mechanics obeys Eq.~\eqref{eq:virial} with the electrostatic
interaction; they are not an independent measurement of it.

\section{Stars}
\textbf{The Sun.} For a star in hydrostatic equilibrium the theorem gives $2K_{\rm th}+U=0$. A contracting star radiates half the gravitational energy it
releases and stores the other half as heat. For a polytrope of index 3, $U=-\tfrac32GM^2/R=-5.69\times10^{41}$\,J for the Sun; contraction alone could
power the present luminosity for $|U|/2L=24$\,Myr. The Sun is 4.6\,Gyr old, so it must burn hydrogen --- the argument of Kelvin and Helmholtz, settled by the
virial theorem.

\textbf{White dwarfs.} In a white dwarf the degenerate electrons become relativistic as the mass grows; the kinetic energy then scales as $1/R$ like the
gravitational energy, and Eq.~\eqref{eq:virial} has a solution only at one mass,
\begin{equation}
M_{\rm Ch}=\frac{\omega_3^0\sqrt{3\pi}}{2}\left(\frac{\hbar c}{G}\right)^{3/2}\frac{1}{(\mu_em_u)^2}=1.456\,M_\odot\quad(\mu_e=2),
\end{equation}
with $\omega_3^0=2.01824$ (Chandrasekhar 1931, 1939). No white dwarf has been found above it; type Ia supernovae mark carbon--oxygen white dwarfs that reach it.

\section{Clusters of galaxies}
Zwicky (1933) applied Eq.~\eqref{eq:virial} to the velocities of the galaxies in the Coma cluster and found far more mass than the light suggests. Virial
(dynamical) masses of clusters are now compared with weak-lensing and X-ray masses in large samples; they agree at the tens-of-percent level, set by
projection, departures from equilibrium in merging systems and non-thermal support.

\textbf{Simulated dark-matter halos.} Cosmological $N$-body simulations measure the virial ratio of every halo directly. Within the virial radius the
published values are $2T/|U|\approx1.1$--$1.3$, rising slowly with mass (Bett et al.\ 2007; Neto et al.\ 2007, ``median slightly greater than unity'';
Power, Knebe \& Knollmann 2012, $\eta=2T/|W|\approx1.15$ at $10^{12}$ to $1.25$ at $10^{15}\,h^{-1}M_\odot$; Klypin et al.\ 2016). The excess over 1 is the
pressure of matter still falling through the boundary: with the surface-pressure term included, Klypin et al.\ find $2K/|W|\approx1.02$--$1.17$. Halos are
called relaxed when $2T/|U|<1.35$ (Neto et al.\ 2007). Collisionless halos, assembled without radiating, sit at the virial equilibrium to within the
infall still under way.

\section{What the kinetic half is}
Let a $1/r$ system fall from rest at infinity ($E_i=0$) into a stable bound state $E_f<0$ and release the difference.
\begin{enumerate}
\item First law: the energy released is $Q=-E_f$.
\item Eq.~\eqref{eq:virial}: $-E_f=\langle K\rangle$, so $Q=\langle K\rangle$.
\item Second law: the transition is irreversible, producing entropy $\Delta S\ge Q/T$ in surroundings at temperature $T$.
\item Landauer (1961): the minimum cost of producing $\Delta S$ is $T\Delta S$; the bound is met when $\Delta S=Q/T$.
\end{enumerate}
\begin{equation}
\langle K\rangle=Q=T\Delta S_{\min}=E_{\rm Landauer}=\tfrac12|\langle V\rangle| .
\end{equation}
Step 2 is the mechanical theorem, so this does not derive Eq.~\eqref{eq:virial}; it identifies its kinetic half with the minimum thermodynamic cost of forming
the bound state. It applies wherever the binding energy is released --- the photon of atomic recombination, the radiation of a contracting cloud or star.
A collisionless system (a dark-matter halo) relaxes without radiating; there the half describes the redistribution of energy in violent relaxation, and the
thermodynamic reading concerns the information written in that relaxation. Jacobson (1995) obtained the Einstein equations from $\delta Q=T\,dS$ on local
horizons; the identity above is the same relation applied at the boundary of a forming bound state.

\section{The black-hole horizon}
A Schwarzschild black hole of mass $M$ has Bekenstein--Hawking entropy and Hawking temperature
\begin{equation}
S=\frac{k_Bc^3A}{4G\hbar}=\frac{4\pi Gk_BM^2}{\hbar c},\qquad T_H=\frac{\hbar c^3}{8\pi Gk_BM}.
\end{equation}
Counting the entropy in bits, $N=S/(k_B\ln2)$, and pricing each bit at Landauer's $k_BT_H\ln2$,
\begin{equation}
N\,k_BT_H\ln2=T_HS=\tfrac12Mc^2 .
\end{equation}
For one solar mass, $N=1.51\times10^{77}$ bits at $T_H=6.17\times10^{-8}$\,K; for the black holes in the Galactic centre ($4.3\times10^6\,M_\odot$) and in
M87 ($6.5\times10^9\,M_\odot$) the ratio is the same, $0.5000000000$. This is the Smarr relation $Mc^2=2T_HS$ of general relativity (Smarr 1973;
Bardeen, Carter \& Hawking 1973). Read with Landauer's principle, half of a black hole's rest energy is the cost of the information on its horizon --- the
same half as in Eq.~\eqref{eq:virial}. For a rotating (Kerr) black hole, $Mc^2=2T_HS+2\Omega_HJ$, and the horizon share is $T_HS/Mc^2=\sqrt{1-\chi^2}/2$:
$0.433$ at $\chi=0.5$, $0.218$ at $\chi=0.9$, $0.032$ at $\chi=0.998$; the rest is rotational energy.

\section{The cosmic horizon}
The Informational Actualization Model (IAM) adds to the horizon first law of Jacobson and of Cai \& Kim (2005) the entropy of the classical records produced as
matter decoheres gravitationally, $-dE=T_H\,d(S_{\rm geo}+S_{\rm info})$. The share of the matter density that enters the record term is the kinetic half of
the binding energy of collapsed structure:
\begin{equation}
\beta_m=\frac{\Omega_m}{2}=0.15765\qquad(\Omega_m=0.3153,\ \text{Planck 2018}).
\end{equation}
The coupling is a prediction. It enters every Planck chain at this value and is never sampled. In the three modified-CAMB chains against the Planck likelihood,
$\Delta\chi^2=+0.54$ against $\Lambda$CDM with the same number of parameters, $\sigma_8$ moves from $0.809$ to $0.800$, and every standard parameter stays
within $0.1\sigma$. Its decisive test is the growth of structure measured by Euclid, where $\mu_0-1=-0.136$ is predicted with $\Sigma_0=0$.

\section{Summary}
\begin{table}[h]\centering\small
\begin{tabular}{llll}
\toprule
System & Scale & Quantity & Value\\\midrule
Hydrogen atom & $10^{-10}$ m & $\langle K\rangle/|\langle V\rangle|$ & $1/2$ ($13.6057$ of $27.2114$\,eV)\\
20 atoms, 10 molecules (HF limit) & $10^{-10}$--$10^{-9}$ m & $T/|V|$ & $1/2$ (exact for converged wavefunctions)\\
White dwarfs & $10^{7}$ m & $M_{\rm Ch}$ & $1.456\,M_\odot$ ($\mu_e=2$)\\
The Sun & $10^{9}$ m & $|U|/2L$ & $24$ Myr (hence nuclear burning)\\
Clusters of galaxies & $10^{23}$ m & dynamical vs lensing mass & agree to tens of percent\\
Simulated halos ($10^{10}$--$10^{15}M_\odot$) & $10^{21}$--$10^{23}$ m & $2K/|W|$ & $1.02$--$1.17$ with surface pressure\\
Schwarzschild black hole & $r_s$ & $T_HS/Mc^2$ & $1/2$ (Smarr)\\
Cosmic horizon & $10^{26}$ m & $\beta_m/\Omega_m$ & $1/2$, predicted; Planck $\Delta\chi^2=+0.54$\\\bottomrule
\end{tabular}
\end{table}

\section*{References}
\small
Bett P.\ et al.\ 2007, MNRAS 376, 215. \quad
Bardeen J.M., Carter B., Hawking S.W.\ 1973, Commun.\ Math.\ Phys.\ 31, 161. \quad
Bekenstein J.D.\ 1973, Phys.\ Rev.\ D 7, 2333. \quad
Bunge C.F., Barrientos J.A., Bunge A.V.\ 1993, At.\ Data Nucl.\ Data Tables 53, 113. \quad
Cai R.-G., Kim S.P.\ 2005, JHEP 02, 050. \quad
Chandrasekhar S.\ 1931, ApJ 74, 81; 1939, \emph{An Introduction to the Study of Stellar Structure}. \quad
Clausius R.\ 1870, Phil.\ Mag.\ 40, 122. \quad
Clementi E., Roetti C.\ 1974, At.\ Data Nucl.\ Data Tables 14, 177. \quad
Hawking S.W.\ 1975, Commun.\ Math.\ Phys.\ 43, 199. \quad
Jacobson T.\ 1995, Phys.\ Rev.\ Lett.\ 75, 1260. \quad
Klypin A.\ et al.\ 2016, MNRAS 457, 4340. \quad
Landauer R.\ 1961, IBM J.\ Res.\ Dev.\ 5, 183. \quad
Neto A.F.\ et al.\ 2007, MNRAS 381, 1450. \quad
Power C., Knebe A., Knollmann S.R.\ 2012, MNRAS 419, 1576. \quad
Planck Collaboration 2020, A\&A 641, A6. \quad
Smarr L.\ 1973, Phys.\ Rev.\ Lett.\ 30, 71. \quad
Zwicky F.\ 1933, Helv.\ Phys.\ Acta 6, 110.
\end{document}
